% from ega1-2-fr.tex
\oldpage[I]{104}
Si $X$ est un $S$-préschéma, $v:X'\to X$ un morphisme de préschémas, le
morphisme composé $X'\xrightarrow{v}X\to S$ définit $X'$ comme
$S$-préschéma~; en particulier, tout préschéma induit sur un ouvert $U$ de
$X$ peut être considéré comme un $S$-préschéma au moyen de l'injection
canonique.

Si $u:X\to Y$ est un $S$-morphisme de $S$-préschémas, la restriction de $u$
à tout préschéma induit sur un ouvert $U$ de $X$ est donc un $S$-morphisme
$U\to Y$. Inversement, soit $(U_\alpha)$ un recouvrement ouvert de $X$ et pour
chaque $\alpha$, soit $u_\alpha:U_\alpha\to Y$ un $S$-morphisme~; si, pour tout
couple d'indices $(\alpha,\beta)$, les restrictions de $u_\alpha$ et de
$u_\beta$ à $U_\alpha\cap U_\beta$ coïncident, il existe un $S$-morphisme
$X\to Y$ et un seul dont la restriction à chaque $U_\alpha$ soit $u_\alpha$.

Si $u:X\to Y$ est un $S$-morphisme tel que $u(X)\subset V$, où $V$ est une
partie ouverte de $Y$, alors $u$, considéré comme morphisme de $X$ dans $V$,
est encore un $S$-morphisme.
\end{env}

\begin{env}[2.5.4]
\label{I.2.5.4-fr}
Soit $S'\to S$ un morphisme de préschémas~; pour tout $S'$-préschéma, le
morphisme composé $X\to S'\to S$ définit $X$ comme $S$-préschéma.
Inversement, supposons que $S'$ soit le préschéma induit sur un ouvert de
$S$~; soit $X$ un $S$-préschéma et supposons que le morphisme structural
$f:X\to S$ soit tel que l'on ait $f(X)\subset S'$~; alors on peut considérer
$X$ comme un $S'$-préschéma. Dans ce dernier cas, si $Y$ est un $S$-préschéma
dont le morphisme structural applique aussi l'espace sous-jacent dans $S'$,
tout $S$-morphisme de $X$ dans $Y$ est aussi un $S'$-morphisme.
\end{env}

\begin{env}[2.5.5]
\label{I.2.5.5-fr}
Si $X$ est un $S$-morphisme, $\varphi:X\to S$ le morphisme structural, on
appelle $S$-\emph{section de $X$} un $S$-morphisme de $S$ dans $X$,
c'est-à-dire un morphisme de préschémas $\psi:S\to X$ tel que
$\varphi\circ\psi$ soit l'identité de $S$. On désigne par $\Gamma(X/S)$
l'ensemble des $S$-sections de $X$.
\end{env}
